\section{Axiomatic conditions for a Galois theory}
\label{V.4}
Let $\cal{C}$ be a category, $F$ a covariant functor from
$\cal{C}$ into the category of finite sets. Suppose the
following conditions are satisfied:

\begin{enumerateb}

\item[(G 1)] $\cal{C}$ has a final object\footnote{Recall that an object
$e$ of~$\cal{C}$ is called a \emph{final object} if for every $X$ in
$\cal{C}$, $\Hom(X,e)$ has exactly one element. One defines
dually an \emph{initial object} of~$\cal{C}$.} and the fiber
product of two objects over a third one in $\cal{C}$
exists (this axiom can also be stated by saying that in
$\cal{C}$ finite projective limits exist).

\item[(G 2)] Finite sums exist in $\cal{C}$ (hence so does an
initial object $\emptyset_{\cal{C}}$, playing the role of the empty
set), as does the quotient of an object of~$\cal{C}$ by a finite
group of automorphisms.

\item[(G 3)] Let $u \colon X \to Y$ be a morphism in $\cal{C}$;
then $u$ factors as a product $X\lto{u'} Y'\lto{u''} Y$, with $u'$ a \emph{strict} epimorphism and
$u''$ a monomorphism which is an isomorphism onto a direct summand
of~$Y$.

\item[(G 4)]
% original p. 119
The functor $F$ is left exact (\ie transforms a right unit
into a right unit, and commutes with fiber
products).

\item[(G 5)] $F$ commutes with finite direct sums, transforms
strict epimorphisms into epimorphisms, and commutes with
passage to the quotient by a finite group of automorphisms.

\item[(G 6)] Let $u \colon X \to Y$ be a morphism in $\cal{C}$ such
that $F(u)$ is an isomorphism; then $u$ is an isomorphism.
\end{enumerateb}

Our object is to construct a topological group~$\pi$, a projective
limit of finite groups, and an equivalence of the
category $\cal{C}$ with the category $\cal{C}(\pi)$
\emph{of finite sets on which} $\pi$ \emph{operates
continuously} (\ie in such a way that the stabilizer of a
point is an open subgroup, or equivalently that there exists a discrete
quotient group which already operates on the set
under consideration), the equivalence constructed transforming the given
functor $F$ into the obvious inclusion functor of~$\cal{C}(\pi)$
into the category of finite sets. Let us note right away that
the category $\cal{C}(\pi)$ constructed by means of a topological
group~$\pi$, and the preceding inclusion functor,
do indeed satisfy conditions (G 1) to (G 6).

We proceed in several steps.
\begin{enumerateb}
\item[a)] \emph{Let} $u \colon X \to Y$ \emph{be in}
$\cal{C}$. \emph{For} $u$ \emph{to be a monomorphism, it is necessary and
sufficient that} $F(u)$ \emph{be one}. (Uses (G 1), (G 4), (G 6)).

Indeed, to say that $u$ is a monomorphism means that the projection
$\pr_1 \colon X \times_Y X \to X$ is an isomorphism.

\item[b)] \emph{Every object} $X$ \emph{of} $\cal{C}$ \emph{is
artinian}.

Indeed, if $X' \to X'' \to X$ are monomorphisms such that
$F(X')$ and $F(X'')$ have the same image in $F(X)$, then by a)
$F(X') \to F(X'')$ is an isomorphism, hence $X' \to X''$ is an
isomorphism by (G 6).

\item[c)] \emph{The functor} $F$ \emph{is strictly
pro-representable}. (\cf Grothendieck, Technique de descente et
th\'eor\`emes d'existence en G\'eom\'etrie Alg\'ebrique, II,
S\'eminaire Bourbaki 195, February 1960).

Indeed, by \loccit prop\ptbl \Ref{V.3.1}, this follows from
b) and (G 4). One can therefore find a projective system over a
filtered ordered set~$I$:
$$
P=(P_i)_{i \in I}
$$
% original p. 120
in $\cal{C}$, considered as a pro-object of~$\cal{C}$, and a
functorial isomorphism
\begin{equation*}
\label{eq:V.4.*} \tag{$*$} F(X)=\Hom_{\Pro(\cal{C})}(P,X)=\varinjlim_{i} \Hom_{\cal{C}}(P_i,X)
\end{equation*}
This isomorphism is realized by an element
$$
\varphi \in \varprojlim_{i} F(P_i)=F(P)
$$
One may moreover assume that the transition homomorphisms
$\varphi_{ji} \colon P_i \to P_j$ ($i \geq j$) are
\emph{epimorphisms}, and that \emph{every epimorphism} $P_i \to
P'$ is equivalent to an epimorphism $P_i \to P_j$ for
suitable $j \leq i$ (which determines the projective system~$P$ in an essentially unique way).

An object $X \in \cal{C}$ is said to be \emph{connected}
if it is not isomorphic to the sum of two objects of~$\cal{C}$
not isomorphic to the initial object $\emptyset_{\cal{C}}$.

\item[d)] \emph{The} $P_i$ \emph{are connected and not isomorphic
to} $\emptyset_{\cal{C}}$.

If $X$ is a left unit, one has $F(X)=\emptyset$ by (G 5)
applied to the direct sum of an empty family, and
conversely by (G 6). Hence if $X'$ is an object of~$\cal{C}$
which is not a left unit, \ie such that $F(X') \neq
\emptyset$, there exists no morphism from~$X'$ into $X$. Hence if some
$P_i$ is a left unit, then $i$ is a largest
element of the filtered ordered index set $I$, and
formula \eqref{eq:V.4.*} would mean $F(X)=\Hom(P_i,X)={}$\kern2pt a set
reduced to one element for every $X$, which is absurd
since $F(\emptyset_{\cal{C}})=\emptyset$. Hence the~$P_i$ are not
isomorphic to $\emptyset_{\cal{C}}$.

Suppose that one has $P_i=A \amalg B$, whence by (G 5)
$F(P_i)=F(A) \amalg F(B)$; in particular the element $a_i$
of~$F(P_i)$, corresponding via \eqref{eq:V.4.*} to
the identity homomorphism $P_i \to P_i$, lies in $F(A) \amalg
F(B)$, for instance in $F(A)$. This means that there exists a $j
\geq i$ such that $\varphi_{ij} \colon P_j \to P_i$ factors as
$P_j \to A \to P_i=A \amalg B$, where the second arrow
is the canonical morphism. Hence $F(P_j) \to F(P_i)$ factors
as $F(P_j) \to F(A) \to F(P_i)=F(A) \amalg F(B)$, and since $F(P_j)
\to F(P_i)$ is surjective by (G 5), it follows that
$F(B)=\emptyset$, hence $B$ is isomorphic to $\emptyset_{\cal{C}}$.

\item[e)]
% original p. 121
\emph{Every morphism} $u \colon X \to Y$ \emph{in} $\cal{C}$,
\emph{with} $X$ \emph{not isomorphic to} $\emptyset_{\cal{C}}$
\emph{and} $Y$ \emph{connected, is a strict epimorphism. Every
endomorphism of a connected object is an automorphism}.

Consider the factorization (G 3) of~$u$; since $X \neq
\emptyset_{\cal{C}}$ it follows from (G 6) that $F(X) \neq \emptyset$, hence
$F(Y') \neq \emptyset$, hence $Y' \neq \emptyset_{\cal{C}}$, hence, $Y$
being connected, $Y'$ is identified with~$Y$, hence $u$ is a
strict epimorphism. Suppose that $u$ is an endomorphism of
the connected object $X$; let us prove that it is an automorphism. Indeed, we
may assume~$X$ not isomorphic to $\emptyset_{\cal{C}}$, hence $u$ is
a strict epimorphism by the foregoing, hence $F(u)$ is
an epimorphism by (G 5), and since $F(X)$ is a finite set,
$F(u)$ is bijective. Hence $u$ is an automorphism by (G 6).

In particular, \emph{every endomorphism of a} $P_i$ \emph{is an
automorphism}.

\item[f)] \emph{The following conditions on a} $P_i$ \emph{are
equivalent}: (i) The natural injective map
$\Hom(P_i,P_i) \to \Hom(P,P_i) \simeq F(P_i)$ is also surjective,
\ie for every $u \colon P \to P_i$ there exists a $v \colon P_i \to
P_i$ such that $u=v\varphi_i$ (where $\varphi_i$ is the canonical
homomorphism $P \to P_i$). (ii) The group $\Aut(P_i)$ operates transitively on~$F(P_i)$. (iii)~The group $\Aut(P_i)$ operates
simply transitively on~$F(P_i)$.

Indeed, identifying $\Hom(P,P_i)$ with $F(P_i)$, the map
considered in (i) is none other than $v \mto
F(v)(\varphi_i)$. The equivalence of the three conditions then follows
from the fact that $\Hom(P_i,P_i)=\Aut(P_i)$ and that the preceding
map is already injective.

A $P_i$ satisfying the equivalent conditions
(i), (ii), (iii)
of~f) is called \emph{Galois}.

\item[g)] \emph{For every} $X$ \emph{in} $\cal{C}$, \emph{there exists
a} Galois $P_i$ \emph{such that every} $u \in \Hom(P,X)$ \emph{factors
as} $P \lto{\varphi_i} P_i \longrightarrow X$.

Let $J=\Hom(P,X)=F(X)$; it is a finite set, hence there exists a
$P_j$ such that every $u \colon P \to X$ factors as $P\lto{j}P_j\longrightarrow X$, or equivalently such that the natural morphism
$$
P \to X^J \qquad (J=\Hom(P,X))
$$
factors as
$$
P \lto{\varphi_j} P_j \to X^J
$$
By
% original p. 122
virtue of (G~3), the morphism $P_j \to X^J$ factors as the product
of a monomorphism and a strict epimorphism, which one may
take of the form $\varphi_{ij} \colon P_j \to P_i$. We are therefore
reduced to proving that $P_i$ is Galois. Let $k$ be an index
$\geq j$ such that every morphism $P \to P_i$ factors through
$P \lto{\varphi_k} P_k \longrightarrow P_i$. Note that the
natural morphism $P_k \to X^J$ again factors as the composite
$$
P_k \lto{\varphi_{ik}} P_i \lto{U} X^J
$$
where the first arrow is a strict epimorphism
by~e), and the second a monomorphism. We want to prove that for
a given morphism $\psi \colon P_k \to P_i$, there exists an
endomorphism $v$ of~$P_i$ such that $\psi = v \varphi_{ik}$. But for
every $u \in \Hom(P_i,X)$, consider $u \psi \in \Hom(P_k,X)$;
it is therefore of the form $u'\varphi_{ik}$ with a well-determined
$u' \in \Hom(P_i,X)$. The map $u \mto u'$ from~$J$ into $J$
thus defined by $\psi$ is moreover injective, since $\psi$ is
an epimorphism by virtue of e); it is therefore bijective since
the set~$J$ is finite. The bijective map $u \mto u'$ from~$J$
into $J$ therefore defines an isomorphism $\alpha \colon X^J \isomto
X^J$ making the diagram
$$
\xymatrix{P_k \ar[r]^-{\varphi_{ik}} \ar@{=}[d] & P_i \ar[r]^-U & X^J
\ar[d]^-{\alpha}_-{\simeq} \\ P_k \ar[r]^{\psi} & P_i \ar[r]^-U & X^J}
$$
commutative. By the uniqueness properties of the factorization
of a morphism as a product of a monomorphism and a strict
epimorphism, it follows (since $\psi$ too is a strict
epimorphism by e)) that one can find a morphism $v \colon P_i \to
P_i$ which leaves the diagram commutative, qed.

We conclude in particular that \emph{the Galois} $P_i$
\emph{form a cofinal system in the system of the} $(P_j)$. We
shall therefore have, since for a Galois object $P_i$ one has
$$
\Hom(P,P_i)=\Hom(P_i,P_i)=\Aut(P_i),
$$
by passing to the limit:
$$
\Hom(P,P) = \varprojlim_{i} \Hom(P,P_i) = \varprojlim_{i} \Hom(P_i,P_i) =
\varprojlim_{i} \Aut(P_i)
$$
where
% original p. 123
the projective limit is taken over the Galois $P_i$. Moreover,
by means of the identification $\Hom(P,P_i)=F(P_i)$, and taking into account that
$F$ transforms epimorphisms into epimorphisms, one sees that the
transition homomorphisms in the preceding projective system
are surjective. We conclude from all this:

\item[h)] \emph{One has}
$$
\Hom(P,P) = \Aut(P) = \varprojlim_{i} F(P_i) = \varprojlim_{i} \Aut(P_i),
$$
\emph{where the projective limit is taken over the Galois}
$P_i$.

In particular, $\Aut(P)$ appears as the projective limit of a
projective system of finite groups (the transition homomorphisms
being surjective); we shall endow it with the projective limit
topology of the discrete topologies. \emph{We shall denote
by}~$\pi$
\emph{and call the fundamental group}
(of~$\cal{C}$ equipped with~$F$) \emph{the group opposite
to}~$\Aut(P)$. This group therefore operates \emph{on the right}
on~$P$; it is the projective limit of finite groups $\pi_i$
operating on the right on the Galois $P_i$, $\pi_i$ being
the group opposite to $\Aut(P_i)$.

Taking into account the functorial isomorphism
$$
F(X)=\Hom(P,X)
$$
and the definition of~$\pi$, one therefore sees that $\pi$ operates
\emph{on the left} on~$F(X)$, and moreover continuously,
by g) (for with the notation of g), it is in fact $\pi_i$
which operates on~$F(X)$). It is trivial that for every morphism $u
\colon X \to Y$ in $\cal{C}$, the morphism $F(u) \colon F(X) \to
F(Y)$ is compatible with the operations of~$\pi$. \emph{One may
therefore henceforth regard} $F$ \emph{as a covariant
functor}
$$
F \colon \cal{C} \to \cal{C}(\pi)
$$
\emph{where} $\cal{C}(\pi)$ \emph{is the category of finite
sets on which} $\pi$ \emph{operates continuously on the left}.

We shall now define a functor in the opposite direction:
$$
G \colon \cal{C} \from \cal{C}(\pi)
$$
by the formula
$$
G(E)=P \times_{\pi} E,
$$
where $P \times_{\pi} E$ is defined as the solution of the
universal problem summarized by
$$
\Hom_{\cal{C}}(P \times_{\pi} E,X) \isomto \Hom_{\pi}(E,\Hom(P,X))
$$
(where
% original p. 124
in the right-hand side $\Hom(P,X)=F(X)$ is regarded
as a set on which $\pi$ operates on the left). One must prove
the existence of the object $P \times_{\pi} E$.

\item[i)] \emph{Let} $Q$ \emph{be an object of} $\cal{C}$ \emph{on which a
finite group} $G$ \emph{operates on the right, and} $E$ \emph{a
finite set on which} $G$ \emph{operates on the left. Then} $G
\times_G E$ \emph{exists, and the canonical map}
$$
F(Q) \times_G E \to F(Q \times_G E)
$$
\emph{is an isomorphism}.

Since finite direct sums exist in $\cal{C}$ by (G 2), and
$F$ commutes with them by (G 5), one is immediately reduced to the case
where $G$ operates transitively on~$E$, for if the~$E_j$ are the
orbits of~$G$ in $E$, one will have
$$
Q \times_G E = \underset {j}{\amalg}\ Q \times_G E_j.
$$
Let then $a \in E$, and let $H$ be its stabilizer; one sees immediately
from the definition that $Q \times_G E$ is identified with $Q/H$.
Whence the existence, thanks to (G~2), and the commutation property
for~$F$, thanks to (G~5).

\item[j)] \emph{Let} $E$ \emph{be an object of} $\cal{C}(\pi)$, \emph{and
let} $P_i$ \emph{be Galois such that} $\pi_i$ \emph{already
operates on} $E$. \emph{Then} $P_i \times_{\pi_i} E$
\emph{exists and one has a canonical isomorphism}
$$
E \isomto F(P_i \times_{\pi_i} E)
$$
\emph{If} $j \geq i$ \emph{is such that} $P_j$ \emph{is Galois,
then the canonical homomorphism} $P_j \times_{\pi_j} E \to\allowbreak P_i
\times_{\pi_i} E$ \emph{is an isomorphism}.

The first assertion is a special case of i), taking into account
that $\pi_i$ operates simply transitively on~$F(P_i)$,
which is equipped with a marked point $\varphi_i$, whence an isomorphism
$F(P_i) \times_{\pi_i} E \simeq E$. For the second assertion one
uses for instance (G 6).

For every $j$, let $\cal{C}_j$ be the full subcategory of
$\cal{C}$ formed by the $X$ such that $\Hom(P_j,X) \to \Hom(P,X)
\simeq F(X)$ is bijective. We know by g) that $\cal{C}$ \emph{is
the filtered union of the} $\cal{C}_j$. We therefore have, for $X \in
\cal{C}_j$:
\begin{align*}
\Hom_{\pi}(E,\Hom(P,X)) & \simeq \Hom_{\pi}(E,\Hom(P_j,X)) \simeq
\Hom_{\pi_j}(E,\Hom(P_j,X))\\ & \simeq \Hom(P_j \times_{\pi_j} E,X)
\end{align*}
and,
% original p. 125
taking into account the last assertion in j), we find an
isomorphism functorial in the object $X$ of~$\cal{C}_j$:
$$
\Hom_{\pi}(E,\Hom(P,X)) \simeq \Hom(P_i \times_{\pi_i} E,X)
$$
Since this is true for every $j$ and these functorial isomorphisms,
for variable $j$, induce one another, we conclude:

\item[k)] \emph{Under the conditions of} j), \emph{the composite
of the canonical morphisms}
$$
E \to \Hom(P_i,P_i \times_{\pi_i} E) \to \Hom(P,P_i \times_{\pi_i} E)
$$
\emph{makes} $P_i \times_{\pi_i} E$ \emph{a solution of the
universal problem defining} $P \times_{\pi} E$,
\ie \emph{the latter exists and one has an isomorphism}
$$
P \times_{\pi} E \isomto P_i \times_{\pi_i} E
$$

This completes the construction of the functor $G(E)$. On the other hand, one has
a functorial homomorphism
$$
\alpha \colon \id_{\cal{C}(\pi)} \to FG
$$
\ie a homomorphism functorial in the object $E$ of~$\cal{C}(\pi)$:
$$
\alpha(E) \colon E \to FG(E) = F(P \times_{\pi} E)
$$
namely the composite of the canonical morphisms
$$
E \to F(P) \times_{\pi} E \to F(P \times_{\pi} E)
$$
(where the first comes from the marked point $\varphi \in F(P)$).
Combining j) and k), one finds:

\item[l)] \emph{The homomorphism} $\alpha$ \emph{is an isomorphism}

One defines likewise a functorial homomorphism
$$
\beta \colon GF \to \id_{\cal{C}}
$$
\ie a homomorphism functorial in the object $X$ of~$\cal{C}$:
$$
\beta(X) \colon P \times_{\pi} F(X) \to X
$$
as
% original p. 126
the one associated with the $\pi$-homomorphism
$$
F(X) \to \Hom(P,X)
$$
inverse to the canonical isomorphism $\Hom(P,X) \isomto F(X)$.

\item[m)] \emph{The composites}
\begin{gather*}
F(X) \lto{\alpha(F(X))} FGF(X) \lto{F(\beta(X))} F(X)\\
G(E) \lto{G(\alpha(E))} GFG(E) \lto{\beta(G(E))} G(E)
\end{gather*}
\emph{are the identity isomorphisms}.

A trotting donkey.

Taking l) into account, it follows:

\item[n)] \emph{The homomorphism} $\beta$ \emph{is an isomorphism}

We have thus obtained the promised result:
\end{enumerateb}

\begin{theorem}
\label{V.4.1}
Let $\cal{C}$ be a category satisfying conditions \emph{(G
1)}, \emph{(G 2)}, \emph{(G~3)} at the beginning of this no., and $F$ a
covariant functor from~$\cal{C}$ into the category of finite
sets, satisfying conditions \emph{(G 4)}, \emph{(G 5)} and
\emph{(G 6)}. Then the preceding canonical constructions
define equivalences of categories $F\colon
\cal{C}\to \cal{C}(\pi)$ and $G\colon \cal{C}(\pi)\to \cal{C}$
quasi-inverse to one another. More precisely, there exists
a pro-object~$P$ of~$\cal{C}$, and a functorial isomorphism $F(X)
\isomfrom \Hom(P, X)$; $\pi$ is the group opposite to the group of
automorphisms of~$P$, topologized in a suitable way,
so that~$\pi$ operates continuously on the
sets $\Hom(P, X)\simeq F(X)$. Finally $G$ is given by
$G(E)\simeq P\times_{\pi} E$.
\end{theorem}

\begin{remarks}
\label{V.4.2}
The statement of conditions (G 1) to (G 6) becomes simpler
and more pleasant if one replaces (G 2) and (G 5) respectively by:
\begin{enumerateb}
\item[(G' 2)] Finite inductive limits exist in
$\cal{C}$.
\item[(G' 5)] The functor $F$ is right exact (\ie commutes with
finite inductive limits).
\end{enumerateb}
These conditions are apparently stronger than (G~2) and (G~5), but
it follows immediately from the structure
theorem~\Ref{V.4.1} that they are implied by (G~1)
to (G~6). Note, however, that in the cases which will
interest us, the verification of (G~2) and (G~5) seems
in fact simpler than that of (G'~2) and (G'~5). I do not know whether,
in condition (G~3), the fact that $u''$ is an isomorphism onto a
direct summand of~$Y$ could not be omitted.
\end{remarks}
